% Generated draft: requires mathematical and source review.
\documentclass[11pt]{article}
\usepackage[margin=25mm]{geometry}
\usepackage{fontspec}
\setmainfont{Latin Modern Roman}
\usepackage{amsmath,amssymb,mathtools,unicode-math}
\setmathfont{Latin Modern Math}
\usepackage{enumitem,needspace,xurl,hyperref}
\setlength{\parindent}{0pt}
\setlength{\parskip}{5pt}
\setlength{\emergencystretch}{3em}
\newcommand{\N}{\mathbb{N}}
\newcommand{\Z}{\mathbb{Z}}
\newcommand{\Q}{\mathbb{Q}}
\newcommand{\R}{\mathbb{R}}
\newcommand{\C}{\mathbb{C}}
\newcommand{\F}{\mathbb{F}}
\newcommand{\A}{\mathbb{A}}
\newcommand{\PP}{\mathbb{P}}
\newcommand{\E}{\mathbb{E}}
\newcommand{\eps}{\varepsilon}
\newcommand{\dd}{\,\mathrm{d}}
\begin{document}
\section*{UnsolvedMath ID: 30006458}
\textbf{Source problem code:} catalogue entry
\subsection{Definitions and mathematical statement}
The supplied record for permanent UnsolvedMath ID \(30006458\) is marked \texttt{statement\_status: unrecoverable}. Its verification note identifies the excerpt as exposition surrounding a reported result, rather than a separately recoverable open problem. Consequently, the following records the displayed stability schema without promoting it to a fully specified conjecture.

Write \(\beta_n\) for the objects appearing in the excerpt, identified as braid groups by the supplied title. The excerpt describes coefficient summands \(V_\lambda\), indexed by irreducible representations of a symplectic group, occurring in a sheaf denoted \(L_k\). It says that their multiplicities depend on \(k\) and \(n\), whereas the family of summands is independent of these parameters. The symplectic groups, coefficient field, sheaves, representations, and their dependence on \(n\) are not defined in the supplied material.

The proposed form of a \emph{uniform stability bound} is
\[
H_k(\beta_n;V_\lambda)\longrightarrow H_k(\beta_{n+1};V_\lambda)
\quad\text{is an isomorphism whenever}\quad k<An-B.
\]
Here the arrows are the stabilization maps referred to in the excerpt; their construction and compatible coefficient maps are not supplied. The word ``uniform'' suggests constants independent of \(\lambda\), but the excerpt does not explicitly quantify \(A,B\), state their allowed values or dependencies, or specify the admissible indices and boundary cases. The displayed formula therefore remains a schema requiring these hypotheses, not a complete quantified assertion.

The excerpt attributes an earlier stability range to a theorem of Randal-Williams--Wahl, displaying its inequality in the extracted text as \(k<21(n-|\lambda|)\). The factor rendered as \(21\) is unresolved and may reflect damaged mathematical typography; neither its intended value nor the meaning of \(|\lambda|\) is supplied reliably enough to formalize this range.

The source reports a calculation of stable homology with a Galois action and says that its topological calculations, together with a suitable uniform stability bound and control of unstable Betti numbers, would imply its ``Conjecture 3.'' It associates the error term with Frobenius traces on unstable homology. Neither Conjecture 3, its moment asymptotic, the necessary estimates, nor the precise Galois action is provided. In particular, the extracted notation \(\operatorname{Gal}(\mathbb Q/\mathbb Q)\) does not specify the intended nontrivial Galois group. No complete statement of that conjecture or implication can be recovered from this record.
\subsection{Short English statement}
The excerpt describes a stabilization isomorphism for twisted braid-group homology in a range \(k<An-B\), intended to be uniform in the representation index. Essential definitions and quantifiers are missing. The supplied verification classifies this as exposition of a reported result, with no recoverable separate open problem; the underlying ``Conjecture 3'' is absent.
\subsection{Sources}
\begin{itemize}
\item Supplied local UnsolvedMath registry record: permanent ID \(30006458\), exact source problem code \texttt{catalogue entry}, title ``Reported Uniform Stability Bound for Twisted Braid-Group Homology.'' Repository registry locations are \texttt{data/unsolvedmath/problems.json} and \texttt{catalog/problems.json}; these files were not independently inspected for this draft.
\item Supplied source citation: \emph{Analytic Number Theory (2026)}, Oberwolfach Reports, DOI 10.4171/owr/2025/51, pp.\ 2701--2764. \url{https://doi.org/10.4171/owr/2025/51}. The background describes Oberwolfach Reports as documenting mathematical workshops and their open-problem discussions. The record gives proposed year 2025 and no proposer.
\item Supplied literature source: ``Uniform twisted homological stability.'' \url{https://arxiv.org/abs/2402.00354}. Authors and a complete bibliographic citation are not supplied.
\item Supplied literature source: ``Oberwolfach Report 51/2025.'' \url{https://doi.org/10.4171/owr/2025/51}.
\item Literature assessment dated exactly \texttt{2026-08-22}: the supplied curation labels the record partially solved and reports that uniform homological stability in a linear range yields the predicted function-field moment asymptotics in a large-\(q\) regime, while the full conjecture for all odd prime powers is not thereby settled. This is supplied dated curation, not an independently verified literature assessment.
\item The supplied statement-verification note says that the official report explains that the required uniform bound is proved in the cited work, and that this extraction is surrounding proof exposition rather than a separate open problem. The record retains the excerpt for provenance and promotes no clean statement. These are supplied verification claims; the linked sources have not been visited for this draft.
\item The excerpt attributes a general stability theorem to Randal-Williams--Wahl under its reference label [4], and discusses work under reference label [2]. Their bibliography entries and the association of [2] with the supplied literature links are not provided.
\end{itemize}
\end{document}
